\documentclass[a4paper]{article}
\usepackage[margin=2.5cm]{geometry} % article's default margins: ~4 cm
\usepackage{titling}
\setlength{\droptitle}{-2cm} % \maketitle leaves a lot of space above the title
\usepackage{amsmath}
\usepackage[md]{juliatex} % [md]: also analysis.md (pandoc)

\title{Scientific Data Analysis}
\author{Miguel Ferrera}

\begin{document}
\maketitle

\section{The model}

We analyze a damped trigonometric function, defined for each new data point.
\begin{equation}
  f(x) = e^{-0.1x} \sin(x).
\end{equation}

\section{Simulation in Julia}

We define the function and evaluate a few points:

\begin{juliacode}
f(x) = exp(-0.1x) * sin(x)
x = 0:0.1:50
println("Points evaluated: ", length(x))
round.(f.(0:2:8), digits = 3)
\end{juliacode}

The plot is generated separately, without showing the code behind it:

\begin{juliarun}
using Plots
plot(x, f.(x), label = "f(x)", lw = 2, xlabel = "Time", ylabel = "Amplitude")
mkpath("resources")
savefig("resources/damped.pdf")
savefig("resources/damped.svg") # for analysis.md (GitHub does not show PDF)
\end{juliarun}

\begin{figure}[htbp]
  \centering
  \juliagraphics[width=0.75\linewidth]{resources/damped.pdf}
  \caption{Damped function generated by Julia at compile time.}%
  \label{fig:damped}
\end{figure}

As figure~\ref{fig:damped} shows, the amplitude decays exponentially.
The computed maximum is \juliav{round(maximum(f.(x)), digits = 4)}.

\end{document}
